\scanpage{5}
mod $p$ des valeurspropres de l'endomorphisme donné de $E(Y)$). Enfin, (et c'est un petit luxe supplémentaire moins important pour nous, car seul les trois derniers termes de (8.1.3) sont pour nous vraiment importants du point de vue assymptotique \ldots) \deleted{\unsure{on aura $\operatorname{Ker}u$}} le \unsure{premier} terme $\operatorname{Ker}u_0$ dans (8.1.3) est nul:
\[
(8.3.3)\qquad\operatorname{Ker}u_0=0.
\]

8.4. Notons maintenant que l'on a des applications semi-linéaires Frob opérant fonctoriellement sur tous les espaces de cohomologie \unsure{envisagés}, et en particulier sur \deleted{\unsure{ceux}}\inserted{les termes} intervenant dans (8.1.3), qui sont donc des vectoriels de dimension finie sur $k$ avec endomorphisme $p$-linéaire. Considérons, pour un tel vectoriel $V$, le plus grand sous-espace \unsure{vectoriel} stable $V_{\mathrm{SS}}$ (partie ``semi-simple'' de $V$) stable par $F_V$, et sur lequel $F_V$ soit \inserted{``semi-simple'' i.e.} surjectif (NB $k$ est parfait); \deleted{\unsure{Le foncteur}} il revient au,même de dire que $F$ est surjectif sur ce sous-espace, et \uline{nilpotent}\deleted{\unsure{r}} sur le quotient. Pour $V$ variable, la formation de $V_{\mathrm{SS}}$ est un foncteur exact en $V$. Ainsi, la suite exacte (8.1.3) nous donne, compte tenu de (8.3.3), la suite exacte d'espace vectoriels à opérateurs $p$-linéaires semi-simples:
\diagram{$(8.4.1)$\quad\begin{tikzcd}[column sep=small,row sep=large]
&&\inserted{({}_p\operatorname{Tors}^n(X))_{\mathrm{ss}}}&&&|[draw,ellipse]|\inserted{0\ \text{en vertu de 8.4.4.}}\arrow[lll,bend right=15]\arrow[d,phantom,"\rotatebox{90}{$=$}"]&\\
0\arrow[r]&(\operatorname{Ker}u_1)_{\mathrm{SS}}\arrow[r]&(\operatorname{Ker}u_2)_{\mathrm{SS}}\arrow[u,phantom,"="]\arrow[r]&(E(Y)\otimes_W k)_{\mathrm{SS}}\arrow[r]&E_{\mathrm{DR}}(Y)_{\mathrm{SS}}\arrow[r]&(\operatorname{Coker}u_2)_{\mathrm{SS}}\arrow[r]&0
\end{tikzcd}.}

\inserted{8.4.2.} D'après la théorie de structure des vectoriels à opération $p$-linéaire semi-simple, \deleted{qui}\inserted{ceux-ci} peuvent s'interpréter aussi comme des Modules libres \unsure{de} type fini sur le faisceau d'anneaux étale constant $\underline F_p$ (corps premier à $p$ éléments), ou encore \deleted{\unsure{comme}} (une fois choisie une clôture séparable) comme des \deleted{\unsure{$G$-$\underline F_p$}} représentations de $G=\operatorname{Gal}(\bar k/k)$ sur des vectoriels de \unsure{dimension} finie sur $\underline F_p$. Dans cette catégorie il y a une autodualité, qui \unsure{s'interprète} indifféremment en termes de dualité d'espaces de dimension finie sur $k$ avec opération $p$-linéaire semi-simple, ou en termes de dualité de

\mspage{4}
